\documentclass{midl} % Include author names

% The following packages will be automatically loaded:
% jmlr, amsmath, amssymb, natbib, graphicx, url, algorithm2e
% ifoddpage, relsize and probably more
% make sure they are installed with your latex distribution

\usepackage{mwe} % to get dummy images
\jmlrvolume{-- Under Review}
\jmlryear{2027}
\jmlrworkshop{Full Paper -- MIDL 2027 submission}
\editors{Under Review for MIDL 2027}

\title[Short Title]{Full Title of Article}

 % Use \Name{Author Name} to specify the name.
 % If the surname contains spaces, enclose the surname
 % in braces, e.g. \Name{John {Smith Jones}} similarly
 % if the name has a "von" part, e.g \Name{Jane {de Winter}}.
 % If the first letter in the forenames is a diacritic
 % enclose the diacritic in braces, e.g. \Name{{\'E}louise Smith}

 % Two authors with the same address
 % \midlauthor{\Name{Author Name1} \Email{abc@sample.edu}\and
 %  \Name{Author Name2} \Email{xyz@sample.edu}\\
 %  \addr Address}

 % Three or more authors with the same address:
 % \midlauthor{\Name{Author Name1} \Email{an1@sample.edu}\\
 %  \Name{Author Name2} \Email{an2@sample.edu}\\
 %  \Name{Author Name3} \Email{an3@sample.edu}\\
 %  \addr Address}


% Authors with different addresses:
% \midlauthor{\Name{Author Name1} \Email{abc@sample.edu}\\
% \addr Address 1
% \AND
% \Name{Author Name2} \Email{xyz@sample.edu}\\
% \addr Address 2
% }

%\footnotetext[1]{Contributed equally}

% More complicate cases, e.g. with dual affiliations and joint authorship
\midlauthor{\Name{Author Name1\midljointauthortext{Contributed equally}\nametag{$^{1,2}$}} \orcid{1111-2222-3333-4444} \Email{abc@sample.edu}\\
\addr $^{1}$ Address 1 \\
\addr $^{2}$ Address 2 \AND
\Name{Author Name2\midlotherjointauthor\nametag{$^{1}$}} \Email{xyz@sample.edu}\\
\Name{Author Name3\nametag{$^{2}$}} \Email{alphabeta@example.edu}\\
\Name{Author Name4\midljointauthortext{Contributed equally}\nametag{$^{3}$}} \Email{uvw@foo.ac.uk}\\
\addr $^{3}$ Address 3 \AND
\Name{Author Name5\midlotherjointauthor\nametag{$^{4}$}} \Email{fgh@bar.com}\\
\addr $^{4}$ Address 4
}

\begin{document}

\maketitle

\begin{abstract}
This is a great paper and it has a concise abstract.
\end{abstract}

\begin{keywords}
List of keywords, comma separated.
\end{keywords}

\section{Introduction}
A fundamental limitation of dermatological AI is that skin generation may not preserve patient-specific skin appearance, or visual patterns that vary with genetics and body site. Healthy skin synthesis models "inpaint" identified lesions with surrounding skin \cite{narges_controllable_2026} or edit disease via diffusion that conforms to an overall learned prior \cite{moutselos_zero-shot_2026}, which doesn’t explicitly preserve the overall patient-specific skin appearance. Skin color disentanglement (or separation of representations) has been done to show lesions across skin tones \cite{yang_automated_2027}. Pathology disentanglement from healthy anatomy in MRI enables synthesizing disease-free scans with pathology removed \cite{xia_pseudo-healthy_2020}. Combining the two ideas, I propose a model that disentangles skin from disease-related change to enable healthy skin synthesis and skin tone modification that preserves patient-specific skin. \\
Chronic Graft-Versus-Host Disease (cGVHD) is the leading cause of late non-relapse mortality after allogeneic stem cell transplantation [4]. Body surface area (BSA) of active epidural disease determines treatment response [4], yet clinician estimates of BSA differ by over 20\% [4]. Photographs of a patient's unaffected skin improve dermatologist agreement [5], but healthy references of patient skin are rarely available, motivating reliable healthy-skin synthesis. cGVHD segmentation models perform worse on darker skin [4], motivating learning representations to bolster model performance.

\section{Related Work}
Models to discuss:
cgDDI -> why do we not like it?
FairDITA -> 
Automated disentangling of lesions  -> we borrow their technique of randomized decolorazation for texture disentanglement, but the idea of disentanglement by compression and having a pigmentation vector is helpful 

Mention HACA3 inductive bias for disentanglement 
Xia et-al is helpful way of thinking of 


\section{Our Framework and Data} 
\subsection{BioSkin-generated Disease}
For "ground truth" healthy counterparts 

Uszing CL ideas from HACA3: (disease to healthy skin representations, pigmentation CL to encourage capturing healthier skin tones (although issue remains with tone regulation for healthy genrations) 

\subsection{Model Architecture}
The network uses three encoders to learn three disentangled representations: a CNN that gives a skin map that is mean to capture high frequency texture details, a CNN that gives a pigmentation vector, and a disease map that captures disease related change in relation to the baseline.\\

\textbf{NOT SO SURE ABOUT RAND IMPLEMENTATION. TRY RMMD (AS ORIGINALLY DONE) AND IF BETTER JUST DO THAT AND DISCARD }
Prior to the patch being passed into the skin encoder, it undergoes randomized decolorization motivated by the approach in \cite{yang_automated_2027}, where there is a monotonic pixel-wise transformation from $[0, 1]^3  \rightarrow [0,1]$. For each training image, non-negative channel weights are sampled from a symmetric Dirichlet distribution and used to form a random convex projection of the RGB channels. The resulting single-channel image is then subjected to random perturbations to destroy intensity cues while preserving spatial organization. This discourages the texture encoder from relying on chromatic and intensity cues when constructing the spatial skin representation. Instead, the skin encoder receives a monochromatic image input, which effectively disentangles it from color-related formation contained in the pigmentation vector. The skin color and appearance is encoded by a 16-dimensional pigmentation vector, which is disentangled from disease-related appearance via \eqref{eq:pig_contrastive}. Together, the pigmentation and skin representations are passed to the decoder to reconstruct a baseline representation of skin without disease, and this output is enforced by an L1 loss to the healthy pixels outside the lesion map (if applicable) \\

Afterwards, the disease encoder's representation adds an RGB residual to the baseline representation that captures the disease-related change pertaining to the patch (if any). The disease representation is disentangled from healthy appearance via \eqref{eq:dis_contrastive}. This is the final reconstruction output that has an L1 loss to the overall input. 



\subsubsection{Contrastive Learning}
In order to encourage the pigemntation vector to be disentangled from disease-related coloration and to encourage it to capture patient-specific pigmentation that relates to healthy skin,  we use a supervised contrastive loss as detailed below. The set of valid anchors is $i \in I \equiv {1...N}$ and $A(i) \equiv I \setminus {i}$. This encourages the pigmentation vector to capture similar skin tone characteristics associated with healthy patches regardless of the disease status of the input patch. \\
\begin{equation}[pig]
\mathcal{L}_{\mathrm{pig}}
=
-\sum_{i\in\mathcal{I}}
\log
\frac{
\exp\left(\mathbf{p}_i^\top \mathbf{p}_p / \tau\right)
}{
\sum_{a\in\mathcal{A}(i)}
\exp\left(\mathbf{p}_i^\top \mathbf{p}_a / \tau\right)
},
\label{eq:pig_contrastive}
\end{equation}

\noindent where
\[
\begin{aligned}
\mathbf{p}_i &: \text{pigmentation encoding of input image } i,\\
\mathbf{p}_p &: \text{pigmentation encoding of healthy patch from same patient and body part},\\
\mathbf{p}_a &: \text{different-patient sample with minimum ITA difference of 45 degrees} \\
\tau &: \text{temperature parameter}.
\end{aligned}
\]
Similarly, to disentangle the disease representation from the skin representation, we use a disease contrastive loss to pull disease representations
with similar disease coverage together while separating them from disease representations extracted from healthy
skin. \\
\begin{equation}
\mathcal{L}_{\mathrm{dis}}
=
-\sum_{i\in\mathcal{I}}
\log
\frac{
\exp\left(\mathbf{d}_i^\top \mathbf{d}_p / \tau\right)
}{
\sum_{a\in\mathcal{A}(i)}
\exp\left(\mathbf{d}_i^\top \mathbf{d}_a / \tau\right)
},
\label{eq:dis_contrastive}
\end{equation}
\noindent where
\[
\begin{aligned}
\mathbf{d}_i &: \text{disease encoding of diseased input image } i,\\
\mathbf{d}_p &: \text{disease encoding of a diseased patch with coverage within 20\% of input},\\
\mathbf{d}_a &: \text{disease encoding of a healthy patch},\\
\tau &: \text{temperature parameter}.
\end{aligned}
\]



\subsubsection{Overall Loss Formulation}


\section{Results}
\subsection{Disease Removal Generation Metrics}



\begin{table}[htbp]
 % The first argument is the label.
 % The caption goes in the second argument, and the table contents
 % go in the third argument.
\floatconts
  {tab:example}%
  {\caption{BioSkin dataset disease removal comparison}}%
  {\begin{tabular}{ll}
  \bfseries Metric & \bfseries Our Model \bfseries cgDDI \\
  LPIPS inside lesion & 0.12345\\
  L1 inside Lesion & 0.67890\\
  Gram-Matrix inside Lesion & 0.54321\\
  LPIPS entire image & 0.09876 \\
  L1 entire image & 0.67890\\
  Gram-Matrix entire image & 0.54321\\
  DICE on input-output diff & \\ 
  
  \end{tabular}}
\end{table}
INSERT IMAGE OF COMPARISON OF DISEASE REMOVAL RESULTS BETWEEN TWO MODELS

\subsection{cGVHD Disease Removal Results}
Limited evals


\subsection{cGVHD Pigmentation and Texture Swap Evaluation}
Train ResNet classification model to predict diseased or non diseased, trained on different augmented training sets, especially for darker skin.
Compare our augmented data vs SCCI \cite{yang_automated_2027}  paper (TRAIN FROM SCRATCH) 
- one without dataset aug (already have code for that)
- eval performance with darker skin (ITA brackets)
- try a naive method too for skin color transfer 

Metrics: LIPIS/texture preservation, pigmentation swap RGB error, CIELAB space degree difference

UNLIKE 

\subsection{Skin Texture Transfer}
LPIPS and Gram-Matrix



\begin{figure}[htbp]
 % Caption and label go in the first argument and the figure contents
 % go in the second argument
\floatconts
  {fig:example}
  {\caption{Example Image}}
  {\includegraphics[width=0.5\linewidth]{example-image}}
\end{figure}

\begin{algorithm2e}
\caption{Computing Net Activation}
\label{alg:net}
 % older versions of algorithm2e have \dontprintsemicolon instead
 % of the following:
 %\DontPrintSemicolon
 % older versions of algorithm2e have \linesnumbered instead of the
 % following:
 %\LinesNumbered
\KwIn{$x_1, \ldots, x_n, w_1, \ldots, w_n$}
\KwOut{$y$, the net activation}
$y\leftarrow 0$\;
\For{$i\leftarrow 1$ \KwTo $n$}{
  $y \leftarrow y + w_i*x_i$\;
}
\end{algorithm2e}

\section{Conclusions}
INCLUDE LIMITATOINS ( numerous, BIoskin not perfect)



\clearpage  % Acknowledgements, references, and appendix do not count toward the page limit (if any)
% Acknowledgments---Will not appear in anonymized version
\midlacknowledgments{We thank a bunch of people.}



\bibliography{midl-samplebibliography}


\appendix



\section{Proof of Theorem 1}

This is a boring technical proof of
\begin{equation}\label{eq:example}
\cos^2\theta + \sin^2\theta \equiv 1.
\end{equation}

\section{Proof of Theorem 2}

This is a complete version of a proof sketched in the main text.

\end{document}
